\subsection{Backward compensation and concentration}

Fix an integer time $T\le T_\zeta$ and a root. For the terminal value $\sum_eZ^\oc_{e,\varnothing}$ we solve
the equation
$\partial_t\Phi+\cL_t\Phi=0$ backward in continuous time $t\in[0,T]$, in the linear space spanned
by~\eqref{eq:patrones}, writing
$\Phi_t=\sum_{e,J}\{c^\oc_{e,J}(t)Z^\oc_{e,J}+c^\di_{e,J}(t)Z^\di_{e,J}\}$. The solution has a closed form.
With $\Delta y=(T-t)/N$
and $z=z_t$, let \[
 \theta_e=\frac{\pi_e\,\Delta y}{z^4},\qquad \xi_{e,i}=\frac{p_{d_i}\,\Delta y}{z};
\] then \begin{equation}\label{eq:compensador_directo}
 \Phi_t=\sum_{e\in\mathcal I_{rc}}\bigl(O_e+\theta_e\chi_e\bigr)
 \prod_{i=1}^4\bigl(V_{d_i}-\xi_{e,i}\chi_{d_i}\bigr),
\end{equation} that is, $c^\oc_{e,J}=\prod_{i\in J}(-\xi_{e,i})$ and
$c^\di_{e,J}=\theta_e\prod_{i\in J}(-\xi_{e,i})$. Indeed, at
$t=T$ the quantities $\theta_e$ and $\xi_{e,i}$ vanish, leaving the terminal value. Since
$\partial_t\Delta y=\partial_tz=-1/N$, we
have $\partial_t\xi_{e,i}=-\beta_i+\alpha\xi_{e,i}$ and $\partial_t\theta_e=-k_e+4\alpha\theta_e$;
differentiating the products, \[
\begin{aligned}
 \partial_tc^\oc_{e,J}&=|J|\,\alpha\,c^\oc_{e,J}+\sum_{i\in J}\beta_i\,c^\oc_{e,J\setminus\{i\}},\\
 \partial_tc^\di_{e,J}&=(4+|J|)\,\alpha\,c^\di_{e,J}+\sum_{i\in J}\beta_i\,c^\di_{e,J\setminus\{i\}}
 -k_e\,c^\oc_{e,J},
\end{aligned}
\] which are the backward equations: by~\eqref{eq:generador}, they make the coefficient of each pattern in
$\partial_t\Phi+\cL\Phi$
vanish. This is the same structure as in Lemma~\ref{lem:compensador}. Since $0\le\Delta y/z\le1$,
$p_{d_i}\le1$, $\pi_e\le C/n$, and
$z\ge\zeta$, we obtain the uniform bounds \begin{equation}\label{eq:coeficientes}
 |c^\oc|\le C_\zeta,\quad |\partial_tc^\oc|\le C_\zeta/n,\quad |c^\di|\le C_\zeta/n,\quad
 |\partial_tc^\di|\le C_\zeta/n^2,
\end{equation} and also $|\partial_t^2c^\oc|=O_\zeta(n^{-2})$ and $|\partial_t^2c^\di|=O_\zeta(n^{-3})$. The
coefficients solve the
continuous-time equation and are evaluated at integer times; the error of this discretization is bounded in
the proof below.

\begin{lema}[Concentration of absorbers]\label{lem:conc_directa} For every fixed tolerance $\gamma>0$,
choosing $\eps$ sufficiently
small in Lemma~\ref{lem:regularidad} gives \begin{equation}\label{eq:conc_directa}
 \Pr\bigl\{|A_T^{rc}-\Phi_0|>\gamma n,\ \sigma'>T\bigr\}\le e^{-c_{\zeta,\gamma}n},
\end{equation} where $\sigma'=\inf\{t:E_t\ge\eps n\text{ or failure}\}$ satisfies
$\Pr(\sigma'\le T_\zeta)\le
e^{-c_{\zeta,\eps}n}$. A union bound over the $n^2$ roots and $O(n)$ times preserves the tail estimate.
\end{lema}

\begin{proof} \emph{Drift.} Let $\Delta c=c(t+1)-c(t)$ for each coefficient. By the backward equation,
$\sum c\,\cL Z=-\sum\partial_tc\,Z$, and therefore \[
 \E(\Phi_{t+1}-\Phi_t\mid\mathcal F_t)=\sum c\,\bigl[\E(\Delta Z\mid\mathcal F_t)-\cL Z\bigr]
 +\sum(\Delta c-\partial_tc)\,Z+\E\Bigl(\sum\Delta c\,\Delta Z\Bigm|\mathcal F_t\Bigr),
\] with the sums over the patterns of the root and the coefficients evaluated at $t$. Since $\sum_eO_e\le n$
and there are $n^2$
squares, Lemma~\ref{lem:generador} and~\eqref{eq:coeficientes} bound the first term by
$C_\zeta(\eps+n^{-1})$. The other two terms
form the discretization error: by the bounds on $\partial_t^2c$, the second is $O_\zeta(n^{-1})$; since
$\sum|\Delta Z^\oc|=O(1)$
and $\sum|\Delta Z^\di|=O(n)$ per root, so is the third. Hence
$|\E(\Phi_{t+1}-\Phi_t\mid\mathcal F_t)|\le C_\zeta(\eps+n^{-1})$.

\emph{Increments.} A destination diagonal affects at most two occupied candidates, since each row and column
contains at most one queen.
Each vertex occurs in $O(n)$ available patterns for this root, with coefficients $O(1/n)$. A birth adds an
occupied term with bounded
coefficient. Accounting also for the change in the coefficients gives $|\Delta\Phi|\le C_\zeta$, including
the exit step.

\emph{Martingale.} Lemma~\ref{lem:regularidad} with tolerance $\eps/2$ bounds $\Pr(\sigma'\le T_\zeta)$. We
freeze
$\widetilde\Phi_s=\Phi_{s\wedge\sigma'}(\text{state}_{s\wedge\sigma'})$, including its time argument. The
drift bound also holds on the
exit step. The martingale
$\widetilde\Phi_s-\widetilde\Phi_0-\sum_{u<s}\E(\Delta\widetilde\Phi_u\mid\mathcal F_u)$ has bounded
increments, and the accumulated drift error is $C_\zeta(2n\eps+2)$. We choose $\eps$ to absorb this error
into the tolerance $\gamma$
before taking $n$ large. Azuma's inequality then gives~\eqref{eq:conc_directa}, since $\Phi_T=A_T^{rc}$.
\end{proof}

\begin{lema}[Initial value]\label{lem:valor_inicial} With $y=T/N$, \begin{equation}\label{eq:valor_inicial}
 \Phi_0=a_n^{rc}(y):=y\sum_{e\in\mathcal I_{rc}}\pi_e\prod_{i=1}^4\bigl(1-y\,p_{d_i}\bigr).
\end{equation} \end{lema}

\begin{proof} Initially, $O_e=0$ and $\chi_v=V_v=1$ on all lines, and at $t=0$ we have $\Delta y=y$ and
$z=1$;
by~\eqref{eq:compensador_directo}, $\Phi_0=\sum_e\theta_e\prod_i(1-\xi_{e,i})$, which
is~\eqref{eq:valor_inicial}. The same follows
by duality. The forward solution of the ideal linear system~\eqref{eq:generador} with these initial data is,
at time $y$ and with
$z=1-y$, \[
 Z^{\di,*}_{e,J}=z^{4+|J|}\prod_{i\notin J}(1-yp_{d_i}),\qquad
 Z^{\oc,*}_{e,J}=y\,\pi_e\,z^{|J|}\prod_{i\notin J}(1-yp_{d_i}),
\] as can be checked by differentiation. The source $k_eZ^{\di,*}_{e,J}$ cancels $z^4$ and the remaining
terms are exactly those
of~\eqref{eq:generador}. Let $Z^*(t)$ be the vector of these values with $y=t/N$, and $c(t)$ the vector of
coefficients of
$\Phi_t$. Then $\frac{d}{dt}\langle c(t),Z^*(t)\rangle=0$, because $c$ solves the backward equation and
$Z^*$ the forward equation
with the same operator; hence
$\Phi_0=\langle c(0),Z^*(0)\rangle=\langle c(T),Z^*(T)\rangle=\sum_eZ^{\oc,*}_{e,\varnothing}(T/N)$.
This is the solution of a linear system, not a factorization of the true law. \end{proof}

\subsection{Limiting trajectory}\label{subsec:trayectoria}

Since $x_*\to0$, the target profiles converge to the full-board profile $\nu_1$, with diagonal vacancies
$u_D,u_A\ge\kappa$. In this
subsection the target profile carries the subscript $x_*$, and symbols without subscript denote the
full-board profile. By
Theorem~\ref{thm:perfil} and~\eqref{eq:balance}, $\pi_e=b_{x_*}v_e/n\,(1+O(n^{-1}))$, where $b_{x_*}v_e$ is
the value of
$\nu_{x_*}$ at the center of $e$, and $p_d=1-u_{d,x_*}+O(n^{-1})$, where $u_{d,x_*}$ is the vacancy of $d$
in the target profile; by
Theorem~\ref{thm:perfil}(c), $\nu_{x_*}\to\nu_1$ and $u_{d,x_*}\to u_d$ uniformly. The quadrature estimates
are uniform, and the $O(n)$ exclusions from $G$, the row, and the column contribute $O(1/n)$. Thus,
uniformly in the root,
\begin{equation}\label{eq:trayectoria}
 \frac{a_n^{rc}(y)}n\longrightarrow\tilde a(y;\bar r,\bar c):=
 y\int_{[0,1]^2}\nu_1(i,j)\prod_{d\in\mathcal D(r,c;i,j)}\bigl[1-y(1-u_d)\bigr]\,di\,dj .
\end{equation} In the limit, $u_d$ is the vacancy of the full-board profile. By~\eqref{eq:regularidad},
$|q_t-qz_t|\le\eps n$ except
with probability $e^{-cn}$, so $1-y$ is the limit in probability of the fraction of real pieces still to be
placed. We do not replace the target profile
by the optimal profile of mass $y$. The function $\tilde a$ satisfies
$\partial_y\tilde a=\int\nu_1\prod_d\vartheta_d
-y\int\nu_1\sum_d(1-u_d)\prod_{\ell\ne d}\vartheta_\ell$, with $\vartheta_d=1-y(1-u_d)$ and $\tilde
a(0)=0$ (birth minus loss), and
depends on the position of the root. On the full board, \begin{equation}\label{eq:a_infinito}
 \tilde a(1;r,c)=\int\nu_1\prod_du_d=a(r,c)\ge\kappa^4>0,
\end{equation} with $a$ as in Section~\ref{sec:intro}. This is a limit of the deterministic profile and
requires no regularity of the
process up to its last step.
